\documentclass[11pt,a4paper,oneside]{book}
\usepackage{90-faelles/latex/studiearkiv}
\hypersetup{pdftitle={Matematik – Bind 4: Matematiske metoder}}
\begin{document}
\frontmatter
\begin{titlepage}
\centering
\vspace*{30mm}
{\sffamily\color{arkivblaa}\Large MATEMATIK TIL UNIVERSITETSFYSIK\par}
\vspace{16mm}
{\Huge\bfseries Bind 4\\[6mm]Matematiske metoder\par}
\vspace{13mm}
{\Large Formler, begreber og matematisk sammenhæng\par}
\vfill
\begin{minipage}{0.8\textwidth}
Kompleks analyse, integraltransformationer, specialfunktioner og beregningsmetoder.

Dette bind genbruger de samme LaTeX-kapitler som det samlede matematikkompendium. Planlagte kapitler er markeret som ufærdige. Der indgår ikke separate gennemregnede opgaveeksempler.
\end{minipage}
\vspace{14mm}
{\small Fysik-studiearkiv · Arbejdsudgave\par}
\end{titlepage}
\chapter*{Læsevejledning}
\addcontentsline{toc}{chapter}{Læsevejledning}
Dette bind er en emneopdelt udgave af det samlede matematikkompendium. Den generelle kilde og alle figurer vedligeholdes centralt under \texttt{00-matematik/chapters} og \texttt{00-matematik/figures}. For hele overblikket samt henvisninger mellem bindene anvendes den samlede udgave. Ufuldstændige kapitler er tydeligt markerede.
\begingroup\small\tableofcontents\endgroup
\mainmatter
\chapter{Komplekse tal og kompleks analyse}
\planlagt{Definitioner, notation, forudsætninger, centrale resultater og fysikforbindelser udvikles samlet.}
\section{Komplekse tal, argument og grene}
\leavevmode\par
% Redaktionelt: tilføj definitioner, gyldighed, relationer og relevante figurer.

\section{Holomorfi og Cauchy–Riemann}
\leavevmode\par
% Redaktionelt: tilføj definitioner, gyldighed, relationer og relevante figurer.

\section{Analytiske og harmoniske funktioner}
\leavevmode\par
% UG-02: analytisk repræsentation, real- og imaginærdel, harmonisk konjugation.

\section{Potensrækker og komplekse Taylorudviklinger}
\leavevmode\par
% UG-03: henvis til emnet Rækker, konvergens og approksimationer for generel konvergensteori.

\section{Konturintegration og Cauchys sætning}
\leavevmode\par
% Redaktionelt: tilføj definitioner, gyldighed, relationer og relevante figurer.

\section{Laurentrækker, singulariteter og residuer}
\leavevmode\par
% Redaktionelt: tilføj definitioner, gyldighed, relationer og relevante figurer.

\section{Konforme afbildninger}
\leavevmode\par
% UG-06: lokal vinkelbevarelse med præcise regularitetsbetingelser.

\chapter{Fourier- og Laplacetransformationer}
\planlagt{Definitioner, notation, forudsætninger, centrale resultater og fysikforbindelser udvikles samlet.}
\section{Konventioner, normalisering og enheder}
\leavevmode\par
% Redaktionelt: tilføj definitioner, gyldighed, relationer og relevante figurer.

\section{Fouriertransformation og inversion}
\leavevmode\par
% Redaktionelt: tilføj definitioner, gyldighed, relationer og relevante figurer.

\section{Laplace og konvergensområder}
\leavevmode\par
% Redaktionelt: tilføj definitioner, gyldighed, relationer og relevante figurer.

\section{Parseval, spektra og differentialoperatorer}
\leavevmode\par
% Redaktionelt: tilføj definitioner, gyldighed, relationer og relevante figurer.

% Kildekobling: UG-07 i pensumregisteret. Udvid disse eksisterende afsnit;
% Fourier-rækker og generel rækkekonvergens henvises til emnet Rækker, konvergens og approksimationer.

\chapter{Specialfunktioner}
\planlagt{Definitioner, notation, forudsætninger, centrale resultater og fysikforbindelser udvikles samlet.}
\section{Gamma- og betafunktionen}
\leavevmode\par
% Redaktionelt: tilføj definitioner, gyldighed, relationer og relevante figurer.

\section{Besselfunktioner og cylindrisk symmetri}
\leavevmode\par
% Redaktionelt: tilføj definitioner, gyldighed, relationer og relevante figurer.

\section{Legendre og sfæriske harmoniske}
\leavevmode\par
% Redaktionelt: tilføj definitioner, gyldighed, relationer og relevante figurer.

\section{Hermite, Laguerre og ortogonalitet}
\leavevmode\par
% Redaktionelt: tilføj definitioner, gyldighed, relationer og relevante figurer.

\chapter{Numeriske metoder}
\planlagt{Definitioner, notation, forudsætninger, centrale resultater og fysikforbindelser udvikles samlet.}
\section{Afrunding, kondition og stabilitet}
\leavevmode\par
% Redaktionelt: tilføj definitioner, gyldighed, relationer og relevante figurer.

\section{Rodfinding og lineære systemer}
\leavevmode\par
% Redaktionelt: tilføj definitioner, gyldighed, relationer og relevante figurer.

\section{Interpolation, differentiation og kvadratur}
\leavevmode\par
% Redaktionelt: tilføj definitioner, gyldighed, relationer og relevante figurer.

\section{ODE-integration, konvergens og fejlestimation}
\leavevmode\par
% Redaktionelt: tilføj definitioner, gyldighed, relationer og relevante figurer.

\end{document}
